\chapter{Tasks Completed}
\label{chap:tasks}

To find the best model for this project, our team followed a strict three-month plan. Here is a breakdown of what we accomplished each month:

\begin{enumerate}
    \item \textbf{Month 1: Research and Setting Up the Basics} 
    
    \begin{itemize}
        \item We started by researching older forecasting methods to understand how inventory prediction used to be done.
        \item We cleaned up the FreshRetailNet-50K dataset and organized it so our models could read it.
        \item We built and trained basic Vanilla RNN and Vanilla LSTM models to see how they performed, which gave us a baseline to beat (44K parameters, 0.785 F1-Score).
    \end{itemize}
    
    \item \textbf{Month 2: Feature Engineering and Building Our Custom Model}
    
    \begin{itemize}
        \item We used Integrated Gradients to figure out which features actually mattered, and threw away the bad ones, leaving us with just 6 strong features.
        \item We designed a new "Dual-Stream" setup that processes crazy demand spikes separately from hard inventory facts.
        \item We finished building our custom \textit{Dual-Stream Inventory Shortcut LSTM}, which shrank the model down to just 4,600 parameters while actually beating the bigger baseline models.
    \end{itemize}
    
    \item \textbf{Month 3: Testing Linear Models and Final Results}
    
    \begin{itemize}
        \item We set up a very strict testing process (using Focal Loss and dynamic threshold scanning) to make sure our results were completely fair and accurate.
        \item We built and tested the \textit{Multi-Kernel DLinear} model to see if it could handle both sudden panic-buying and longer weekly trends.
        \item We collected all our final numbers, created the graphs, proved that the Multi-Kernel DLinear was our best model, and wrote this final thesis report.
    \end{itemize}
\end{enumerate}
